\documentclass[10pt,a4paper]{amsart}
\usepackage[margin=19mm]{geometry}
\usepackage{amsmath,amssymb,mathtools}
\usepackage[scr=rsfs,cal=euler]{mathalfa}
\usepackage{xfrac}
\usepackage[unicode,hidelinks,bookmarks=false]{hyperref}
\newcommand{\ie}{i.e.\ }
\newcommand{\cf}{cf.\ }
\newcommand{\resp}{respectively\ }
\newcommand{\Subsection}{\subsection}
\providecommand{\nobreakdash}{\nobreak\hbox{-}}
\setlength{\emergencystretch}{2em}
\multlinegap0pt
\usepackage{bm}
\usepackage{bbm}
\usepackage{dsfont}
\usepackage{xspace}
\usepackage{cancel}
\usepackage{mathtools}
\usepackage{cleveref}
\usepackage{amsthm}
\makeatletter
\@ifundefined{lemma}{\newtheorem{lemma}{Lemma}}{}
\@ifundefined{proposition}{\newtheorem{proposition}{Proposition}}{}
\makeatother
\DeclareMathOperator{\PGL}{PGL}
\DeclareMathOperator{\PSL}{PSL}
\DeclareMathOperator{\Syl}{Syl}
\DeclarePairedDelimiter{\gen}{\langle}{\rangle}
\title[The structure of finite groups invariably generated by two e]{The structure of finite groups invariably generated by two elements of prime order}
\author{Inna Capdeboscq, Chris Parker}
\date{Source: arXiv:2608.03935v1 (CC BY 4.0)}
\begin{document}
\maketitle

\noindent\emph{Source and licence.} The structure of finite groups invariably generated by two elements of prime order. Version arXiv:2608.03935v1 (CC BY 4.0), distributed under the Creative Commons Attribution 4.0 International licence, as recorded in the arXiv licence metadata for this exact version and shown on the abstract page https://arxiv.org/abs/2608.03935v1. Full rights notice: licenses/arxiv-2608.03935-CC-BY-4.0.txt.\par\smallskip

\noindent\textit{Editorial scope.} Counted unit: the Lemma whose source label is \texttt{lem:no repeated factors} in the article (source0.tex lines 1256--1258, source label \texttt{lem:no repeated factors}), delivered together with the article's own complete original proof of it (source0.tex lines 1260--1280, one \texttt{proof} environment of 21 lines). No mathematics is written, repaired or re-derived: statement and proof are the article's own text line for line. Required context retained uncounted: lem:quotients (lines 394--395); lem:psl2 (lines 905--906). Every definition, display and label of those lines is retained unchanged. Omitted from this package and not redistributed: 1--393, 396--904, 907--1255, 1259--1259, 1281--1306. Nothing inside a retained excerpt is omitted or altered: every retained body is delivered exactly as the article's lines are written, so no omission receipt of any kind accompanies this package. Citations: the retained excerpts carry no citation command at all, so no bibliography entry of the article is transcribed and no reference is redistributed. Reference adaptation: none was needed and none was made; every reference inside the delivered unit resolves inside it, so no unresolved reference or citation is produced. Printed slips kept literal: no printed word, symbol, hypothesis or number of the counted statement or of its proof was altered. Layout and class: the article's own class is \texttt{amsart} with packages including biblatex, amssymb, tikz, enumerate, mathtools, xcolor, textcomp, amsthm, thmtools, hyperref, cleveref; the wrapper is the lane's pinned \texttt{amsart} editorial wrapper at 10pt on A4, which restates the theorem-like environments the retained text needs and adds the article's own macro definitions that the retained text needs (\texttt{\textbackslash PGL}, \texttt{\textbackslash PSL}, \texttt{\textbackslash Syl}, \texttt{\textbackslash gen}), with extra wrapper packages (bm, bbm, dsfont, xspace, cancel, mathtools, cleveref). The delivered numbering is the unit's own. Assets: no retained span contains a figure environment, a graphics-inclusion command, a PDF-page inclusion, a table or a TikZ picture; no image file is copied into this package, the package declares no asset dependency and no image file is redistributed with it. Licence: version arXiv:2608.03935v1 of this article is distributed under the Creative Commons Attribution 4.0 International licence, as recorded in the arXiv licence metadata for this exact version and shown on the abstract page https://arxiv.org/abs/2608.03935v1. A full rights notice is delivered at licenses/arxiv-2608.03935-CC-BY-4.0.txt. Characters: every retained excerpt is pure ASCII, so the delivered engine, XeLaTeX with its default Latin Modern text font, reports no missing character.
\begin{lemma}\label{lem:quotients} Suppose that $G$ is invariably $\Lambda$-generated. Then every quotient of $G$ is invariably $\Lambda$-generated.
\end{lemma}

\begin{proposition}\label{lem:psl2}
 Suppose that $r$ is a prime, $K \cong \PSL_2(r^a)$. Then $G$ is invariably $(2,3)$-generated if and only if $G \cong \PGL_2(3^{2^b})$ for some $b\ge 1$.  \end{proposition}

\begin{lemma}\label{lem:no repeated factors}
Suppose that $H \cong \PGL_2(3^{2^a})$ and   $X= H\times H$. Then no subgroup of $X$ containing $F^*(X)$ is invariably $(2,3)$-generated.
\end{lemma}

\begin{proof} Set $q=3^{2^a}$,   $H_1=\{(h,1)\mid h \in H\}$, $H_2=  \{(1,h)\mid h \in H\}$, $K_i = H_i \cap F^*(X)$ and define $K=K_1K_2=F^*(X)$.
Suppose that $G$ is invariably $(2,3)$-generated with $K \le G \le X$.



 For $1\ne N<K$, a  normal subgroup of $G$, we have $N= K_1$, or $N= K_2$.  Hence $G/N$ is isomorphic to a subgroup of $2 \times \PGL_2(q)$.  By \cref{lem:quotients}, $G/N$ is invariably $(2,3)$-generated as are all its quotients. Since the elementary abelian subgroup of order $2^2$ is not invariably $(2,3)$-generated,  $G \ne X$. By  \cref{lem:psl2}, $\PSL_2(q)$ is not invariably $(2,3)$-generated and so $G/N \not \cong 2 \times \PSL_2(q)$. It follows that $G/N \cong \PGL_2(q)$ and that $|X:G|=2$.

Assume that $(\zeta,\delta)$ is an invariable $(2,3)$-sequence for $G$. Then $\zeta \not \in K$.

   We know $G = K\gen{(z,z)}$ with $z\in H\setminus F^*(H)$ of order $2$. Then $\zeta=(z_1,z_2)$ with $z_1$ and $z_2$ of order $2$ and conjugating by elements of $K_1$ and $K_2$ we see that $\zeta$ is $G$-conjugate to $(z,z)$. Hence we may assume $\zeta=(z,z)$.

      Let $D \in \Syl_3(G)$ contain $\delta$, $D_1= D\cap H_1$ and $D_2=D\cap H_2$. Then $D=D_1D_2$ and $N_G(D)$ has order $q^2(q-1)^2/2$ with $N_G(D)/D$ isomorphic to the direct product of  cyclic groups of orders $q-1$ and $(q-1)/2$. Thus $N_G(D  )$ has four conjugacy classes of elements of order $3$: one class in $D_1$, one class in $D_2$ and two classes on the diagonal of $D_1D_2$ each of length $(q-1)^2/2$.


      If $\delta \in D_i$ with $i=1$ or $2$,  then $\langle D_i,\zeta \rangle \le H_i\gen{\zeta} < G$, a contradiction as $\gen{\zeta,\delta}= G$.   Hence $\delta$ is not in $D_1$ or $D_2$.

      Since $N_G(D)$ controls $3$-fusion in $G$ by Burnside's fusion theorem, $G$ has two conjugacy classes of elements of order $3$  diagonal between $H_1$ and $H_2$. The representatives are $(d,d)$ and $(d,d^z)$ where $d$ in $H$ has order $3$. Indeed, suppose that $g=(az^\ell,bz^\ell) \in G$ with $\ell \in \{0,1\}$ and $a,b \in H'$ is such that $(d,d^z)=(d,d)^g$.
      Then $(d,d^z) = (d^{az^\ell},d^{bz^\ell})$ implies that $d= d^{az^\ell}$ from which we conclude that $z^{\ell}=1$. But then $d^z=d^b$ with $b \in H'$, a contradiction as $d$ and $d^z$ are not conjugate in $H'$.

   Set $L= \{(h,h) \mid h \in H\}< G$ and $M= \{(h,h^z) \mid h \in H\}<G$. Then $L$ and $M$ both contain $\zeta $. Since $(d,d) \in L$ and $(d,d^z) \in M$, we have no candidates for $\delta$. This proves the claim.
\end{proof}

\end{document}
